\subsection{Non-radiating and unreachable modes}
\label{sec:ch471}
% TOC tag: [USE]

Non-radiating currents and unreachable far-zone components lie in \(\ker\mathcal{R}_\omega\) and related synthesis nullities
from Section~\ref{sec:ch27} \cite{stratton1941,harrington2001}. Subsection~\ref{sec:ch471} states their impact on
\(\mathcal{S}_{\mathrm{real}}(\mathcal{C})\): such modes alter near fields without changing \(\mathbf{c}\) \cite{devaney2004}.

\paragraph{Assumptions.}
Import \(\ker\mathcal{R}_\omega\), non-radiating templates Eq.~\eqref{eq:ch2.nr-current-def}, and synthesis null space
Eq.~\eqref{eq:ch4.synthesis-null-space} \cite{stratton1941,harrington2001}. \(\boldsymbol{\Gamma}_\omega\) is compact on
\(\mathcal{J}_{\mathrm{real}}\) \cite{harrington2001,devaney2004}.

\paragraph{Unreachable synthesis directions.}
Define
\begin{equation}
\label{eq:ch4.unreachable-modes}
\mathcal{U}_{\mathrm{unreach}}
=
\ker\boldsymbol{\Gamma}_\omega^{\dagger}
\cap
\mathbb{C}^{I},
\end{equation}
the coefficient directions with zero coupling to every admissible \(\Jimp\in\mathcal{J}_{\mathrm{real}}\)
\cite{stratton1941,harrington2001,devaney2004}. Targets with component along \(\mathcal{U}_{\mathrm{unreach}}\) fail
\(\delta_{\mathrm{syn}}=0\) in Eq.~\eqref{eq:ch4.synthesis-residual-norm} regardless of MoM residual \cite{harrington2001}.

\begin{theorem}[Non-radiating lift does not enlarge \(\mathbf{c}\)]
\label{thm:ch4.nr-lift-invisibility}
For every \(\mathbf{J}_{\mathrm{nr}}\in\ker\mathcal{R}_\omega\cap\mathcal{J}_{\mathrm{real}}\) and every
\(\Jimp\in\mathcal{J}_{\mathrm{real}}\),
\(\boldsymbol{\Gamma}_\omega[\Jimp+\mathbf{J}_{\mathrm{nr}}]=\boldsymbol{\Gamma}_\omega[\Jimp]\) on \(\mathcal{F}_{\mathrm{rad}}\)
\cite{stratton1941,harrington2001}. Near-field data distinguish \(\mathbf{J}_{\mathrm{nr}}\); far-zone \(\mathbf{c}\) does not
\cite{devaney2004}.
\end{theorem}
\begin{proof}
Linearity of \(\mathcal{R}_\omega\) and definition of \(\ker\mathcal{R}_\omega\) \cite{stratton1941}; \(\mathcal{C}\) bounded on
\(\mathcal{F}_{\mathrm{rad}}\) \cite{hansen1988,harrington2001}.
\end{proof}

\paragraph{Worked illustration (NR current superposition).}
Adding a non-radiating current to a feed leaves \(\mathbf{c}\) unchanged while altering reactive near fields
\cite{stratton1941,harrington2001}. Inverse solvers that fit near fields may absorb \(\mathbf{J}_{\mathrm{nr}}\) without changing
\(\mathcal{S}_{\mathrm{real}}(\mathcal{C})\) \cite{devaney2004}.

\paragraph{Problem (unreachable target component).}
\textbf{Problem.} \(\mathbf{c}_{\mathrm{target}}\) has unit weight on a sector nullified by symmetry. Is inversion well-posed?

\textbf{Step 1.} Project \(\mathbf{c}_{\mathrm{target}}\) off \(\mathcal{U}_{\mathrm{unreach}}\) \cite{harrington2001}.

\textbf{Step 2.} Test \(\delta_{\mathrm{syn}}\) on the projection \cite{stratton1941,devaney2004}.

\textbf{Step 3.} Report unreachable component as certification failure, not feed error \cite{hansen1988}.

\textbf{Physical meaning.} Unreachable modes are synthesis null directions, not small coupling \cite{stratton1941}.

\paragraph{Validity.}
Theorem~\ref{thm:ch4.nr-lift-invisibility} assumes linear Maxwell maps on declared domains \cite{stratton1941}. Loss and dispersion
modify \(\ker\mathcal{R}_\omega\) without promoting NR content to \(\mathcal{F}_{\mathrm{rad}}\) \cite{harrington2001}.


\subsection{Reactive subspaces and evanescent accessibility}
\label{sec:ch472}
% TOC tag: [HOME]

Evanescent subspace \(\mathcal{F}_{\mathrm{ev}}\) from Eq.~\eqref{eq:ch2.evanescent-subspace-def} carries reactive storage accessible
to near probes yet invisible to \(\boldsymbol{\Gamma}_\omega\) on \(\mathcal{F}_{\mathrm{rad}}\) \cite{stratton1941,harrington2001}.
Subsection~\ref{sec:ch472} separates reactive accessibility from synthesis accessibility \cite{devaney2004}.

\paragraph{Assumptions.}
Presuppose Proposition~\ref{prop:ch4.evanescent-invisibility} and Subsection~\ref{sec:ch273} reconstructability caps
\cite{stratton1941,harrington2001}. Observation rank \(N\) finite \cite{devaney2004}.

\paragraph{Reactive accessibility functional.}
Define
\begin{equation}
\label{eq:ch4.reactive-accessibility}
\mathcal{A}_{\mathrm{react}}^{(N)}(\mathbf{c})
=
\big\|\mathcal{P}_{\mathrm{react}}\mathbf{c}\big\|_{\mathrm{react}},
\end{equation}
with \(\mathcal{P}_{\mathrm{react}}\) projecting onto reactive components of the observation functional when near data are used
\cite{harrington2001,stratton1941}. Large \(\mathcal{A}_{\mathrm{react}}^{(N)}\) with small \(\|\mathbf{c}\|_{\mathrm{rad}}\) signals
evanescent-dominated certification \cite{devaney2004}.

\begin{proposition}[Reactive accessibility does not enlarge synthesis]}
\label{prop:ch4.reactive-not-synthesis}
\(\mathbf{c}\in\mathcal{S}_{\mathrm{real}}(\mathcal{C})\) is independent of \(\mathcal{A}_{\mathrm{react}}^{(N)}\) on
\(\mathcal{F}_{\mathrm{rad}}\) \cite{stratton1941,harrington2001}. Increasing near-field probe rank \(N\) may raise
\(\mathcal{A}_{\mathrm{react}}^{(N)}\) without enlarging \(\mathcal{S}_{\mathrm{real}}(\mathcal{C})\) \cite{devaney2004,hansen1988}.
\end{proposition}

\paragraph{Worked illustration (near-field scan).}
A scan resolving \(k_{t}>k_{0}\) weight reports high \(\mathcal{A}_{\mathrm{react}}^{(N)}\) while \(\mathbf{c}\) remains
dipole-dominated \cite{harrington2001,jackson1999}. Far-zone OAM claims must use \(\mathcal{K}_{\mathrm{prop}}\) only
\cite{stratton1941,devaney2004}.

\paragraph{Problem (reactive accessibility versus synthesis).}
\textbf{Problem.} Near-field tomography recovers reactive content. Does that enlarge \(\mathcal{S}_{\mathrm{real}}(\mathcal{C})\)?

\textbf{Step 1.} Apply \(\Pi_{\mathrm{rad}}\Pi_{\mathrm{prop}}\) before listing \(\mathbf{c}\) \cite{stratton1941}.

\textbf{Step 2.} Test \(\delta_{\mathrm{syn}}\) on far-zone \(\mathbf{c}\) alone \cite{harrington2001,devaney2004}.

\textbf{Step 3.} Treat reactive recovery as certification, not synthesis enlargement \cite{hansen1988}.

\textbf{Physical meaning.} Reactive accessibility certifies storage; synthesis exports propagating rank \cite{stratton1941}.

\paragraph{Validity.}
Eq.~\eqref{eq:ch4.reactive-accessibility} assumes linear probes and declared calibration \cite{harrington2001}. Loss modifies
\(\mathcal{F}_{\mathrm{ev}}\) split without exporting evanescent weight to \(\mathbf{c}\) \cite{stratton1941}.


\subsection{Synthesis, reconstruction, and inversion limits}
\label{sec:ch473}
% TOC tag: [HOME]

Inverse design and source reconstruction minimize residuals on \(\mathbf{c}\) or near fields; realizability requires target membership in
\(\operatorname{ran}\boldsymbol{\Gamma}_\omega\) \cite{stratton1941,harrington2001,devaney2004}. Subsection~\ref{sec:ch473} states
inversion limits as operator-range theorems \cite{hansen1988}.

\paragraph{Assumptions.}
Presuppose Theorem~\ref{thm:ch4.moore-penrose-synthesis} and Principle~\ref{prin:ch4.operator-chain}
\cite{harrington2001,devaney2004,stratton1941}. Tolerance \(\varepsilon\) fixed \cite{hansen1988}.

\paragraph{Inversion residual decomposition.}
For target \(\mathbf{c}_{\mathrm{target}}\),
\begin{equation}
\label{eq:ch4.inversion-residual-decomposition}
\mathbf{c}_{\mathrm{target}}
=
\mathcal{P}_{\mathrm{ran}}\mathbf{c}_{\mathrm{target}}
+
\mathbf{r}_{\perp},
\qquad
\mathbf{r}_{\perp}\in\ker\boldsymbol{\Gamma}_\omega^{\dagger},
\end{equation}
with \(\mathcal{P}_{\mathrm{ran}}\) from Principle~\ref{prin:ch4.operator-chain} \cite{stratton1941,harrington2001}. Realizability
requires \(\mathbf{r}_{\perp}=\mathbf{0}\); MoM may minimize \(\|\mathbf{r}_{\perp}\|\) only if regularization is undeclared
\cite{devaney2004,hansen1988}.

\begin{theorem}[Inversion limit at rank \(d_{\mathrm{eff}}\)]
\label{thm:ch4.inversion-limit}
Best rank-\(d_{\mathrm{eff}}(\varepsilon)\) approximation of \(\boldsymbol{\Gamma}_\omega\) achieves error
\(\le\varepsilon\sigma_{1}\) in operator norm, and any \(\mathbf{c}_{\mathrm{target}}\) with
\(\|\mathbf{r}_{\perp}\|>\varepsilon\|\mathbf{c}_{\mathrm{target}}\|\) lies outside \(\mathcal{S}_{\mathrm{real}}(\mathcal{C})\)
\cite{stratton1941,harrington2001,devaney2004,hansen1988}.
\end{theorem}
\begin{proof}
Theorem~\ref{thm:ch4.radiative-rank} and Eckart--Young truncation \cite{harrington2001,stratton1941}.
\end{proof}

\paragraph{Worked illustration (Tikhonov masking).}
Tikhonov inversion can suppress \(\mathbf{r}_{\perp}\) in discrete norms while \(\delta_{\mathrm{syn}}>0\) in the continuum gauge
\cite{harrington2001,devaney2004}. \(\delta_{\mathrm{syn}}=0\) is the synthesis certificate \cite{stratton1941}.

\paragraph{Problem (inversion versus realizability).}
\textbf{Problem.} Optimizer reports \(\|\mathbf{T}\mathbf{I}_{s}-\mathbf{c}_{\mathrm{target}}\|<10^{-4}\). Certify realizability.

\textbf{Step 1.} Decompose \(\mathbf{c}_{\mathrm{target}}\) via Eq.~\eqref{eq:ch4.inversion-residual-decomposition} \cite{stratton1941}.

\textbf{Step 2.} Test \(\delta_{\mathrm{syn}}=0\) \cite{harrington2001,devaney2004}.

\textbf{Step 3.} Reject if \(\mathbf{r}_{\perp}\neq\mathbf{0}\) despite small MoM norm \cite{hansen1988}.

\textbf{Physical meaning.} Inversion finds closest coefficients; synthesis demands range membership \cite{stratton1941}.

\paragraph{Validity.}
Theorem~\ref{thm:ch4.inversion-limit} assumes linear forward models \cite{stratton1941}. Nonlinear materials require separate
\(\mathcal{J}_{\mathrm{real}}\) classes \cite{harrington2001}.


\subsection{Localization versus angular resolution}
\label{sec:ch474}
% TOC tag: [HOME][USE SS3.8.4][USE SS3.6.4]

Representation dependence of angular labels in Subsection~\ref{sec:ch384} and localization bounds in
Subsection~\ref{sec:ch364} constrain simultaneous spatial concentration and angular resolution on \(\mathcal{F}_{\mathrm{rad}}\)
\cite{harrington2001,hansen1988,stratton1941}. Subsection~\ref{sec:ch474} states the tradeoff as a realizability inequality on
\(\mathcal{S}_{\mathrm{real}}(\mathcal{C})\) \cite{devaney2004}.

\paragraph{Assumptions.}
Import Eq.~\eqref{eq:ch4.localization-spectrum-product} and Eq.~\eqref{eq:ch3.OAM-representation-dependence}
\cite{stratton1941,hansen1988,harrington2001}. Compact \(\Omega_{s}\) \cite{jackson1999}.

\paragraph{Resolution functional.}
Define
\begin{equation}
\label{eq:ch4.angular-resolution-functional}
\mathcal{R}_{\mathrm{ang}}(\mathbf{c})
=
\sup\left\{
\ell:\ \sum_{m,\sigma}|\widehat{c}_{\ell m}^{(\sigma)}|^{2}
\ge
\eta_{\mathrm{blk}}\|\mathbf{c}\|_{\mathrm{rad}}^{2}
\right\},
\end{equation}
the highest \(\ell\) block carrying fraction \(\eta_{\mathrm{blk}}\) of radiative power \cite{hansen1988,stratton1941}. On \(B_{a}\),
\begin{equation}
\label{eq:ch4.localization-resolution-tradeoff}
\mathcal{L}_{2}[\Jimp]\,\mathcal{R}_{\mathrm{ang}}(\mathbf{c})
\ \gtrsim\
\mathcal{K}_{\mathrm{lr}}(k_{0}a,\eta_{\mathrm{blk}}),
\end{equation}
combining Eq.~\eqref{eq:ch4.localization-spectrum-product} with tail bounds Eq.~\eqref{eq:ch4.high-ell-decay}
\cite{jackson1999,hansen1988,devaney2004}.

\begin{corollary}[Super-resolution obstruction]
\label{cor:ch4.super-resolution-obstruction}
Simultaneous \(\mathcal{L}_{2}\ll a^{2}\) and \(\mathcal{R}_{\mathrm{ang}}\gg k_{0}a\) on the same \(\mathbf{c}\) implies
\(\mathbf{c}\notin\mathcal{S}_{\mathrm{real}}(\mathcal{C})\) on \(B_{a}\) \cite{hansen1988,stratton1941,devaney2004}.
\end{corollary}

\paragraph{Worked illustration (tight beam, small aperture).}
A \(0.25\lambda\) aperture claiming \(2^{\circ}\) beam and \(\ell=8\) OAM violates
Eq.~\eqref{eq:ch4.localization-resolution-tradeoff} \cite{jackson1999,harrington2001,hansen1988}.

\paragraph{Problem (OAM label versus resolution).}
\textbf{Problem.} Report uses LG index \(m_{\mathrm{LG}}=5\) on a paraxial beam. Transfer to VSH realizability?

\textbf{Step 1.} Apply Eq.~\eqref{eq:ch3.paraxial-vs-spherical-OAM} \cite{harrington2001,allen1992}.

\textbf{Step 2.} Test \(\mathcal{R}_{\mathrm{ang}}(\mathbf{c})\) against Eq.~\eqref{eq:ch4.ell-max-scaling} \cite{hansen1988}.

\textbf{Step 3.} Reject essentialist OAM claims that skip basis change \cite{stratton1941,belinfante1940}.

\textbf{Physical meaning.} Angular resolution and localization are jointly bounded by support \cite{hansen1988}.

\paragraph{Validity.}
Eqs.~\eqref{eq:ch4.localization-resolution-tradeoff} assume open-region compact sources \cite{stratton1941}. Cavities quantize
resolution differently via Subsection~\ref{sec:ch372} \cite{harrington2001}.

Section~\ref{sec:ch47} has separated synthesis rank from reactive accessibility, inversion limits, and localization--resolution
tradeoffs \cite{stratton1941,harrington2001,devaney2004}. Section~\ref{sec:ch48} develops scaling laws and spectral concentration
on the same geometry \cite{hansen1988}.


\section{Scaling Laws and Spectral Concentration}
\label{sec:ch48}

Effective dimension \(d_{\mathrm{eff}}\) counts export directions; scaling exponents and concentration inequalities describe how that
count grows with aperture, frequency, and signal bandwidth \cite{stratton1941,jackson1999,hansen1988}. Section~\ref{sec:ch48}
supplies the asymptotic laws and concentration bounds that make truncation inevitable on finite hardware \cite{harrington2001,devaney2004}.
